\documentclass[11pt,a4paper]{article}
\usepackage[italian]{babel}
\usepackage[margin=2.5cm]{geometry}
\usepackage{parskip}
\usepackage{amsmath, amssymb, amsthm}
\usepackage{xcolor}
\usepackage{listings}
\usepackage{tikz}
\usepackage[most]{tcolorbox}
\usepackage{hyperref}

% --- Riquadri con bordo nero, senza numero ---
\tcbset{nota/.style={colback=white, colframe=black, boxrule=0.5pt,
  sharp corners}}
\NewTColorBox{definizione}{}{nota,
  before upper={\textbf{Definizione.}\ }}
\NewTColorBox{osservazione}{}{nota,
  before upper={\textbf{Osservazione.}\ }}
\theoremstyle{definition}
\newtheorem*{esempio}{Esempio}
\newtheorem*{esercizio}{Esercizio}

% --- Codice C ---
% Uso: \begin{codice} ... \end{codice}; \begin{risultato} ... \end{risultato}
%      \lstinline|printf("ciao");| per il codice dentro al testo
\lstdefinestyle{c}{
  language=C,
  basicstyle=\ttfamily\small,
  keywordstyle=\color{blue}\bfseries,
  commentstyle=\color{gray}\itshape,
  stringstyle=\color{red!70!black},
  numbers=left,
  numberstyle=\tiny\color{gray},
  numbersep=8pt,
  columns=flexible,
  keepspaces=true,
  breaklines=true,
  showstringspaces=false,
  tabsize=4,
  morekeywords={include, define},
  literate={à}{{\`a}}1 {è}{{\`e}}1 {é}{{\'e}}1 {ì}{{\`i}}1
           {ò}{{\`o}}1 {ù}{{\`u}}1
}
\lstset{style=c}
\newtcblisting{codice}{listing only, nota, left=6mm,
  listing options={style=c}}
\newtcblisting{risultato}{listing only, colback=black!5, colframe=black,
  boxrule=0.5pt, sharp corners,
  listing options={basicstyle=\ttfamily\small, numbers=none, language={}}}

% --- Diagrammi di flusso (simboli standard) ---
\usetikzlibrary{shapes.geometric, arrows.meta, positioning}
\tikzset{
  terminale/.style = {ellipse, draw, minimum width=2.5cm, minimum height=0.9cm},
  io/.style        = {trapezium, trapezium left angle=70, trapezium right angle=110,
                      draw, minimum width=2.5cm, minimum height=0.9cm},
  processo/.style  = {rectangle, draw, minimum width=2.5cm, minimum height=0.9cm},
  decisione/.style = {diamond, draw, aspect=2, minimum width=2.5cm},
  freccia/.style   = {thick, -{Stealth}}
}

\title{Commenti, identificatori e parole chiave}
\author{Marco Cerbella}
\date{Lezione 6 -- 5 ottobre 2026}

\begin{document}
\maketitle

\section{Commenti}

In C ci sono due tipi di commenti:
\begin{itemize}
    \item \emph{a blocco}: iniziano con \lstinline|/*| e finiscono con \lstinline|*/|, possono
          occupare più righe;
    \item \emph{di riga}: iniziano con \lstinline|//| e finiscono a fine riga.
\end{itemize}

\begin{codice}
double pi = 3.1415926536;   // Pi è 3,14
\end{codice}

All'interno di una stringa o di una costante carattere i simboli \lstinline|/*| e
\lstinline|//| non iniziano un commento:
\begin{codice}
printf("Comments in C begin with /* or //.\n");
\end{codice}

\begin{osservazione}
I commenti a blocco \emph{non si possono annidare}: il primo \lstinline|*/|
chiude il commento aperto. Si possono invece usare per ``spegnere'' del codice
che contiene commenti di riga.
\end{osservazione}

\begin{esempio}
Qui il commento aperto alla riga 3 termina al primo \lstinline|*/| (riga 6): il
secondo \lstinline|*/| è un errore di sintassi.
\begin{codice}
#include <stdio.h>
int main() {
    int x = 4; /*
    int y = 10; /*

    printf("x = %d\n", x); */
    printf("y = %d\n", y); */
}
\end{codice}
\end{esempio}

\subsection{Scrivere buoni commenti}

Il codice si scrive per essere eseguito da un computer ma letto da persone:
da noi stessi mentre lo scriviamo e dopo settimane o mesi, dai colleghi che
devono integrarlo e da chi lo manterrà anni dopo.
\begin{itemize}
    \item Evitare commenti ridondanti, come \lstinline|i = i + 1; // incrementa i|.
    \item Non lasciare codice ``eliminato'' sotto forma di commento (va bene
          solo temporaneamente durante il debug).
    \item Evitare commenti decorati (riquadri di asterischi): sono difficili da
          mantenere.
    \item Un commento che chiarisce spesso segnala codice troppo complesso:
          meglio semplificare il codice, che deve essere \emph{autoesplicativo}.
\end{itemize}

\section{Caratteri}

Il C definisce due insiemi di caratteri: quello \emph{sorgente} (usabile nel
codice) e quello di \emph{esecuzione} (interpretabile dal programma in
esecuzione); nella maggior parte delle implementazioni coincidono. I caratteri
\emph{spazio bianco} sono: spazio, \verb|\f| (form feed), \verb|\n| (a capo),
\verb|\t| (tab orizzontale), \verb|\v| (tab verticale).

\section{Identificatori}

\begin{definizione}
Un \emph{identificatore} è il nome di una variabile, funzione, macro,
struttura o altro oggetto definito in un programma C. Può contenere:
\begin{itemize}
    \item lettere \texttt{a--z} e \texttt{A--Z} (il C è \emph{case-sensitive});
    \item il carattere underscore \texttt{\_};
    \item cifre \texttt{0--9}, ma il primo carattere \emph{non} può essere una cifra.
\end{itemize}
\end{definizione}

\begin{esempio}
Identificatori validi: \texttt{x}, \texttt{dollar}, \texttt{While},
\texttt{error\_handler}, \texttt{scale64}.

Non validi: \texttt{1st\_rank} (inizia con una cifra), \texttt{switch} (parola
riservata), \texttt{y/n} e \texttt{x-ray} (caratteri non ammessi).
\end{esempio}

Non c'è un limite alla lunghezza, ma i compilatori considerano significativi
solo i primi 31 o 63 caratteri. Molti nomi sono già usati dalla libreria
standard (ad esempio \lstinline|printf|) e non si possono riutilizzare per
funzioni proprie o variabili globali.

\section{Parole riservate}

Le seguenti 44 parole chiave sono riservate in C (erano 32 in C89), hanno un
significato preciso per il compilatore e non si possono usare come
identificatori:

\begin{center}
\begin{minipage}{0.9\textwidth}
\ttfamily\small
auto, break, case, char, const, continue, default, do, double, else, enum,
extern, float, for, goto, if, inline, int, long, register, restrict, return,
short, signed, sizeof, static, struct, switch, typedef, union, unsigned, void,
volatile, while, \_Alignas, \_Alignof, \_Atomic, \_Bool, \_Complex, \_Generic,
\_Imaginary, \_Noreturn, \_Static\_assert, \_Thread\_local
\end{minipage}
\end{center}

\subsection{Scegliere i nomi}

Un nome descrive l'oggetto, il suo comportamento e il suo uso: una variabile
chiamata male può confondere molto.
\begin{itemize}
    \item evitare nomi inutilmente lunghi;
    \item usare nomi diversi in ambiti diversi, per evitare confusione.
\end{itemize}

\begin{esempio}
Confrontare \lstinline|a = b * c;| con
\lstinline|weekly_pay = hours_worked * pay_rate;|.
\end{esempio}
\end{document}
